\documentclass[10pt,letterpaper]{article}
\usepackage[T1]{fontenc}
\usepackage[utf8]{inputenc}
\usepackage{lmodern}
\usepackage[margin=0.76in]{geometry}
\usepackage[table]{xcolor}
\usepackage{array,tabularx,longtable,booktabs}
\usepackage{enumitem}
\usepackage{amssymb,amsmath}
\usepackage{hyperref}
\usepackage{fancyhdr,lastpage}
\usepackage{microtype}
\usepackage{parskip}
\definecolor{cdblue}{HTML}{203957}
\definecolor{cdgray}{HTML}{586677}
\hypersetup{colorlinks=true,linkcolor=cdblue,urlcolor=cdblue,pdftitle={Chapter 11: Managing Infrastructure and Environments}}
\pagestyle{fancy}
\fancyhf{}
\lhead{\small\color{cdgray}Continuous Delivery Engineering}
\rhead{\small\color{cdgray}Study guide}
\cfoot{\thepage\ / \pageref{LastPage}}
\setlength{\headheight}{14pt}
\setlength{\emergencystretch}{2em}
\setlist[itemize]{leftmargin=*,itemsep=3pt,topsep=4pt}
\setlist[enumerate]{leftmargin=*,itemsep=3pt,topsep=4pt}
\begin{document}
\begin{center}
{\LARGE\bfseries\color{cdblue}Chapter 11: Managing Infrastructure and Environments}\\[7pt]
{\large III: The Delivery Ecosystem}\\[6pt]
{\small Reference edition: 2026-10-30 $\bullet$ Original educational and operational material}
\end{center}
\noindent\color{cdblue}\rule{\linewidth}{1.2pt}\color{black}
\vspace{0.4em}
\section{Chapter scope and provenance}
Primary reading: Humble and Farley, Continuous Delivery (2010), Chapter 11, "Managing Infrastructure and Environments"; Part III: The Delivery Ecosystem. This guide is an original synthesis and not a reproduction of the chapter. The book provides the foundational delivery principles; current AWS, Terraform and GitHub Actions references in the applied sections are supplementary.


\section{Learning outcomes}
\begin{enumerate}
\item Identify infrastructure ownership and drift detection.
\item Design disposable parity-environments.
\item Define logging and operational observability requirements.
\end{enumerate}

\section{Core conceptual model}
Focus areas: environment automation; provisioning; monitoring; operations collaboration.

\begin{itemize}
\item Input: a version-controlled change, explicit requirements, and a reproducible execution context.
\item Mechanism: controlled feedback, automation, evidence capture, and a stop condition where risk exceeds tolerance.
\item Output: a verifiable release candidate or a diagnosed failure with an accountable owner.
\end{itemize}

\section{Reading and analysis protocol}
\begin{enumerate}
\item Read the chapter and record five propositions, distinguishing an author claim from an example or assumption.
\item Identify two failure modes, one recovery technique, and one measurable acceptance condition.
\item Trace each technique to the stages and controls in the deployment-pipeline-architecture reference.
\item State how the proposed practice changes for a regulated system with segregated duties.
\end{enumerate}

\section{Detailed technical study notes}
\subsection{Environment reproducibility}
Environment templates should describe network, identity, middleware and runtime inputs that influence application behavior.

\subsection{Provisioning ownership}
Differentiate shared platform controls from application-owned resources and audit the IAM handoff between them.

\subsection{State and concurrency}
Infrastructure-as-code requires controlled state storage and locking; concurrent applies can corrupt expected state.

\subsection{Monitoring and alerting}
Emit logs, metrics and traces correlated to deploy IDs. An alert must trigger a known human or automated response.

\subsection{Cloud elasticity}
Ephemeral test environments increase test parallelism but can undermine validity if resource limits and costs are not controlled.

\subsection{Drift response}
Detect manual and automatic changes separately and resolve unauthorized drift rather than accepting it as new baseline.

\section{Domain-specific challenge questions}
\begin{enumerate}
\item For environment reproducibility, identify the precondition, expected evidence, adversarial failure case and measurable corrective action.
\item For provisioning ownership, identify the precondition, expected evidence, adversarial failure case and measurable corrective action.
\item For state and concurrency, identify the precondition, expected evidence, adversarial failure case and measurable corrective action.
\item For monitoring and alerting, identify the precondition, expected evidence, adversarial failure case and measurable corrective action.
\end{enumerate}
\section{Guided exercise}
Plan a Terraform-managed sandbox with remote state isolation and tagging.

\subsection{Required evidence}
\begin{itemize}
\item A version-controlled design or executable artifact.
\item A failure demonstration, negative test or counterexample.
\item An improvement statement with baseline, target, owner and verification method.
\end{itemize}


\section{Engineering scenario}
Staging uses a shared manually patched VM while production uses immutable instances.

\begin{itemize}
\item Describe the failure propagation path and the earliest feasible detection point.
\item Recommend a minimal remediation and its rollback or containment plan.
\item Define a measure that would demonstrate risk reduction.
\end{itemize}

\section{Knowledge checks}
\begin{enumerate}
\item Which condition is the gate, and which action is the response when it fails?
\item What evidence distinguishes successful automation from an assumed successful release?
\item Where might environments, data, dependencies or people invalidate the technique?
\item What feedback could be moved earlier without losing necessary confidence?
\item Which assumptions are book principles, and which are modern supplementary implementation choices?
\end{enumerate}

\section{Completion rubric}
Score each dimension from 0 (missing) to 4 (demonstrated and independently verifiable): conceptual accuracy; reproducibility; negative testing; security and safety; documentation and traceability. Pass threshold: 16/20 with no zero on reproducibility or safety.


\section{Prerequisites and onward route}
Chapters 1--4 and prior dependency stages. Next: Chapters 12 onward.



\end{document}
